% !TEX program = xelatex
% ============================================================================
%   << title >> —— 由「题记」导出于 << now >>
%   编译：xelatex main.tex（跑两遍，目录里的页码才会对上）
%   Overleaf：Menu → Compiler 选 XeLaTeX。
%   想改配色、字体：看下面「配色」「字体」两节，改完重新编译即可。
% ============================================================================
\documentclass[11pt,a4paper]{article}
\usepackage[UTF8,fontset=none,scheme=plain,linespread=1.42]{ctex}
\usepackage[a4paper,top=27mm,bottom=25mm,left=24mm,right=24mm,headsep=9mm,footskip=13mm]{geometry}
\usepackage{amsmath,mathtools}
\usepackage{unicode-math}
\usepackage[table]{xcolor}
\usepackage{graphicx}
\usepackage{array}
\usepackage[most]{tcolorbox}
\usepackage{tikz}
\usetikzlibrary{calc}
\usepackage{fancyhdr}
\usepackage{needspace}
\usepackage[hidelinks,bookmarksopen=true]{hyperref}
\hypersetup{pdftitle={<< title >>},pdfauthor={<< author >>},pdfcreator={题记}}

% ---------------------------------------------------------------- 配色
\definecolor{tjink}{HTML}{141C2E}    % 正文墨色
\definecolor{tjink2}{HTML}{465168}   % 次要文字
\definecolor{tjink3}{HTML}{737D91}   % 注释、页眉
\definecolor{tjrule}{HTML}{D6DCE5}   % 细线、边框
\definecolor{tjgrid}{HTML}{E9EDF3}   % 网格
\definecolor{tjpaper}{HTML}{F3F5F8}  % 封面底色
\definecolor{tjlvl}{HTML}{8196B5}    % 等高线、编号
\definecolor{tjhl}{HTML}{EBF1FA}     % 浅色底
\definecolor{tjaccent}{HTML}{2A78D6} % 强调蓝：想法
\definecolor{tjred}{HTML}{E34948}    % 强调红

% ---------------------------------------------------------------- 字体
% 西文 STIX Two、数学 STIX Two Math、等宽 IBM Plex Mono，都是 TeX Live 自带。
% 中文正文霞鹜新致宋，想法用霞鹜文楷，标题用思源宋体 Heavy；找不到时退回 Fandol 字体。
\setmainfont{STIX Two Text}
\setsansfont{IBM Plex Sans}
\setmonofont{IBM Plex Mono}
\setmathfont{STIX Two Math}
\IfFontExistsTF{LXGW Neo ZhiSong}
  {\setCJKmainfont{LXGW Neo ZhiSong}[AutoFakeBold=2.4]}
  {\setCJKmainfont{FandolSong-Regular.otf}[BoldFont=FandolSong-Bold.otf]}
\IfFontExistsTF{LXGW Neo XiHei}
  {\setCJKsansfont{LXGW Neo XiHei}[AutoFakeBold=2]\setCJKmonofont{LXGW Neo XiHei}}
  {\setCJKsansfont{FandolHei-Regular.otf}[BoldFont=FandolHei-Bold.otf]\setCJKmonofont{FandolHei-Regular.otf}}
\IfFontExistsTF{LXGW WenKai GB Lite}
  {\setCJKfamilyfont{tjkai}{LXGW WenKai GB Lite}[AutoFakeBold=2]}
  {\setCJKfamilyfont{tjkai}{FandolKai-Regular.otf}}
\IfFontExistsTF{Source Han Serif SC Heavy}
  {\setCJKfamilyfont{tjdisplay}{Source Han Serif SC Heavy}}
  {\setCJKfamilyfont{tjdisplay}{FandolSong-Bold.otf}[AutoFakeBold=3]}
\newcommand\tjkai{\CJKfamily{tjkai}}
\newcommand\tjdisplay{\CJKfamily{tjdisplay}\bfseries}

% ---------------------------------------------------------------- 常见的 KaTeX 写法，补上定义免得报错
\providecommand\R{\mathbb{R}}\providecommand\N{\mathbb{N}}\providecommand\Z{\mathbb{Z}}
\providecommand\Q{\mathbb{Q}}\providecommand\C{\mathbb{C}}
\providecommand\bm[1]{\symbf{#1}}
\providecommand\argmin{\operatorname*{arg\,min}}\providecommand\argmax{\operatorname*{arg\,max}}
\providecommand\rank{\operatorname{rank}}\providecommand\tr{\operatorname{tr}}
\providecommand\diag{\operatorname{diag}}\providecommand\sgn{\operatorname{sgn}}
\providecommand\dif{\mathop{}\!\mathrm{d}}

% ---------------------------------------------------------------- 版面
\setlength{\parindent}{0pt}
\setlength{\parskip}{.4em}
\setlength{\emergencystretch}{3em}
\AtBeginDocument{\color{tjink}%
  \setlength{\abovedisplayskip}{6pt plus 2pt}\setlength{\belowdisplayskip}{6pt plus 2pt}%
  \setlength{\abovedisplayshortskip}{3pt}\setlength{\belowdisplayshortskip}{4pt}}
\newcommand\tjmeta{\sffamily\footnotesize\color{tjink3}}
\newcommand\tjdots{\leaders\hbox to .85em{\hss\textcolor{tjrule}{·}\hss}\hfill}
\newcommand\tjdoublerule{\par{\color{tjink}\hrule height 1.1pt}\vspace{1.6pt}{\color{tjink}\hrule height .35pt}}

\pagestyle{fancy}
\fancyhf{}
\fancyhead[L]{\tjmeta\rightmark}
\fancyhead[R]{\tjmeta << title >>}
\fancyfoot[C]{\ttfamily\footnotesize\color{tjink3}\thepage}
\renewcommand\headrulewidth{.4pt}
\renewcommand\headrule{{\color{tjrule}\hrule height\headrulewidth width\headwidth}\vskip-\headrulewidth}

% ---------------------------------------------------------------- 组件
% 关键词标签；本章的关键词用实心
\newtcbox{\tjchip}{on line, enhanced, boxrule=.5pt, arc=6pt, colframe=tjrule, colback=white,
  left=5pt, right=5pt, top=1.2pt, bottom=1.2pt, boxsep=0pt, fontupper=\sffamily\scriptsize\color{tjink2}}
\newtcbox{\tjchipon}{on line, enhanced, boxrule=.5pt, arc=6pt, colframe=tjink, colback=tjink,
  left=5pt, right=5pt, top=1.2pt, bottom=1.2pt, boxsep=0pt, fontupper=\sffamily\scriptsize\color{white}}

% 原题框：左上角贴一个深色「原题」标签；原图放在虚线下方
\newtcolorbox{tjproblem}{enhanced, breakable,
  colback=white, colframe=tjrule, boxrule=.6pt, arc=3pt,
  left=14pt, right=14pt, top=14pt, bottom=10pt, middle=8pt,
  before skip=13pt, after skip=0pt,
  title={原\,题}, fonttitle=\sffamily\scriptsize\bfseries, coltitle=white,
  attach boxed title to top left={xshift=13pt, yshift=-\tcboxedtitleheight/2},
  boxed title style={colback=tjink, colframe=tjink, boxrule=0pt, arc=1.5pt,
    left=5pt, right=5pt, top=1.4pt, bottom=1.4pt},
  segmentation style={tjrule, dashed, line width=.5pt},
  before upper={\setlength{\parskip}{.45em}},
  before lower={\setlength{\parskip}{0pt}}}

\newcommand\tjimage[2]{%
  {\tjmeta\scriptsize 原图\par}\vspace{4pt}%
  {\centering\includegraphics[width=#2\linewidth,height=7.5cm,keepaspectratio]{#1}\par}}

% 想法：左侧一道竖线，文楷。#1 标签（我的想法 / 想法 2 / AI 思路）  #2 颜色
% 自己写的用蓝线，AI 给的思路用灰线、正文略淡，一眼分得清
\newtcolorbox{tjidea}[2]{enhanced, breakable, blanker,
  borderline west={1.8pt}{0pt}{#2},
  left=14pt, right=4pt, top=1pt, bottom=1pt, before skip=11pt, after skip=0pt,
  fontupper=\tjkai\fontsize{11.5}{15}\selectfont,
  before upper={{\sffamily\scriptsize\bfseries\color{#2}#1\par}\vspace{-1pt}\setlength{\parskip}{.35em}}}
\newenvironment{tjaiidea}{\begin{tjidea}{AI 思路}{tjink3}\color{tjink2}}{\end{tjidea}}

% 每道题的抬头：编号、出处、日期、关键词
\newcommand\tjhead[4]{% #1 编号  #2 出处  #3 日期  #4 书签文字
  \par\needspace{9\baselineskip}\vspace{20pt}%
  \pdfbookmark[1]{#4}{tjbm#1}%
  \noindent{\ttfamily\fontsize{9}{9}\selectfont\color{tjlvl}No.}\hspace{1.5pt}%
  {\ttfamily\fontsize{19}{19}\selectfont\color{tjink}#1}\hspace{11pt}%
  {\bfseries\large #2}\hfill{\ttfamily\scriptsize\color{tjink3}#3}\par}
\newcommand\tjkeywords[1]{\vspace{4pt}\noindent#1\par}

% 关键词一章的开头
\newcommand\tjchapter[4]{% #1 序号  #2 关键词  #3 题数  #4 锚点
  \par\needspace{16\baselineskip}\vspace{30pt}%
  \pdfbookmark[0]{#2}{tjkwbm#1}\phantomsection\label{#4}%
  \markboth{#2}{#2}%
  \noindent{\ttfamily\scriptsize\color{tjink3}KEYWORD\,·\,#1}\par\vspace{2pt}%
  \noindent{\tjdisplay\fontsize{25}{32}\selectfont #2}\hfill{\ttfamily\small\color{tjink3}#3\,题}%
  \vspace{6pt}\tjdoublerule\vspace{2pt}}

% 一道题属于多个关键词时，在其余几章里留一条「另见」
\newtcolorbox{tjsee}{enhanced, colback=tjhl, colframe=tjhl, boxrule=0pt, arc=3pt,
  left=11pt, right=11pt, top=7pt, bottom=7pt, before skip=22pt, after skip=0pt,
  before upper={{\sffamily\scriptsize\bfseries\color{tjaccent}本章另见\par}\vspace{4pt}\setlength{\parskip}{3pt}}}
\newcommand\tjseeitem[4]{% #1 编号  #2 出处  #3 摘要  #4 锚点
  \noindent\hyperref[#4]{{\ttfamily\small No.#1}}\hspace{9pt}{\small #2}\hspace{7pt}%
  {\small\color{tjink2}#3}\tjdots{\ttfamily\small\color{tjink3}\pageref{#4}}\par}

% 目录
\newcommand\tjtockw[3]{% #1 关键词  #2 题数  #3 锚点
  \par\needspace{4\baselineskip}\vspace{11pt}%
  \noindent\hyperref[#3]{{\sffamily\bfseries\color{tjink} #1}}\hspace{7pt}{\ttfamily\scriptsize\color{tjink3}#2\,题}%
  \tjdots{\ttfamily\small\pageref{#3}}\par\vspace{1pt}}
\newcommand\tjtocitem[4]{% #1 编号  #2 出处  #3 摘要  #4 锚点
  \noindent\hspace{1.4em}\hyperref[#4]{{\ttfamily\small\color{tjink2}No.#1}}\hspace{9pt}%
  {\small #2}\hspace{7pt}{\small\color{tjink3}#3}\tjdots{\ttfamily\small\color{tjink3}\pageref{#4}}\par}

\begin{document}

% ============================================================== 封面
\begin{titlepage}
\begin{tikzpicture}[remember picture, overlay]
  \fill[tjpaper] (current page.south west) rectangle (current page.north east);
  % 封面小图：f(x,y)=(x^2+4y^2)/2 的等高线，和从 (2,1) 出发的最速下降折线
  \begin{scope}[shift={($(current page.south)+(0,6.9cm)$)}, x=1.45cm, y=1.45cm]
    \filldraw[fill=white, draw=tjrule, line width=.5pt, rounded corners=7pt] (-3.75,-2.05) rectangle (3.75,2.05);
    \begin{scope}
      \clip[rounded corners=7pt] (-3.75,-2.05) rectangle (3.75,2.05);
      \draw[step=.5, tjgrid, line width=.4pt] (-4,-2.5) grid (4,2.5);
      \draw[tjrule!80!tjink3, line width=.45pt] (-4,0) -- (4,0) (0,-2.5) -- (0,2.5);
      \foreach \c in {3.6,2.8284,1.4553,0.7487,0.3852,0.1982}
        \draw[tjlvl, line width=.55pt, fill=tjlvl, fill opacity=.05] (0,0) ellipse ({\c} and {\c/2});
      \draw[tjink3, line width=.5pt, dash pattern=on .8pt off 3pt, line cap=round] (-4,-2) -- (4,2) (-4,2) -- (4,-2);
      \fill[tjink, opacity=.045] (0,0) -- (2,1) -- (1.2,-.6) -- cycle;
      \draw[line width=1.5pt, line cap=round, tjaccent] (2,1) -- (1.4118,-.1765);
      \draw[line width=1.5pt, line cap=round, tjred] (1.4118,-.1765) -- (.5294,.2647);
      \draw[line width=1.5pt, line cap=round, tjaccent] (.5294,.2647) -- (.3737,-.0467);
      \draw[line width=1.5pt, line cap=round, tjred] (.3737,-.0467) -- (.1401,.0701);
      \draw[line width=1.5pt, line cap=round, tjaccent] (.1401,.0701) -- (.0989,-.0124);
      \foreach \p in {(1.4118,-.1765),(.5294,.2647),(.3737,-.0467)}
        \filldraw[fill=white, draw=tjink, line width=.6pt] \p circle (2.3pt);
      \fill[tjink] (2,1) circle (3pt);
      \node[font=\itshape\small, text=tjink, above right=1pt] at (2,1) {$z_0$};
      \node[font=\itshape\small, text=tjink, below right=1pt] at (1.4118,-.1765) {$z_1$};
    \end{scope}
  \end{scope}
  \node[anchor=north, font=\sffamily\scriptsize, text=tjink3] at ($(current page.south)+(0,3.65cm)$)
    {$f(x,y)=\tfrac12(x^2+4y^2)$ 的等高线与最速下降折线：每一步都和上一步垂直};
  \node[anchor=south west, font=\sffamily\footnotesize, text=tjink3] at ($(current page.south west)+(24mm,13mm)$)
    {导出于 << now >><% if author %>\quad·\quad << author >><% endif %>};
  \node[anchor=south east, font=\ttfamily\footnotesize, text=tjink3] at ($(current page.south east)+(-24mm,13mm)$)
    {PROBLEM NOTEBOOK};
\end{tikzpicture}

\vspace*{12mm}
{\ttfamily\footnotesize\color{tjink3}PROBLEM\ \ NOTEBOOK\ \ ·\ \ << year >>}\par
\vspace{10mm}
{\tjdisplay\fontsize{58}{66}\selectfont << title >>}\par
<% if subtitle %>
\vspace{6mm}
{\fontsize{14}{22}\selectfont\color{tjink2}<< subtitle >>}\par
<% endif %>
\vspace{13mm}
{\color{tjrule}\hrule height .5pt}
\vspace{7mm}
\noindent\begin{tabular}{@{}l@{\hspace{17pt}}!{\color{tjrule}\vrule width .5pt}@{\hspace{17pt}}l@{\hspace{17pt}}!{\color{tjrule}\vrule width .5pt}@{\hspace{17pt}}l@{}}
{\tjmeta 收录} & {\tjmeta 关键词} & {\tjmeta 时间} \\[3pt]
{\ttfamily\fontsize{24}{26}\selectfont << count >>}\,{\small 题} &
{\ttfamily\fontsize{24}{26}\selectfont << kw_count >>}\,{\small 个} &
{\ttfamily\fontsize{13}{26}\selectfont << span >>}
\end{tabular}
<% if kw_counts %>

\vspace{9mm}
\begingroup\setlength{\parskip}{5pt}\noindent
<% for k, c in kw_counts %>
\tjchip{<< k >>\,\textcolor{tjink3}{<< c >>}}\hspace{3pt}%
<% endfor %>
\par\endgroup
<% endif %>
\end{titlepage}

% ============================================================== 目录
\markboth{目录}{目录}
\pdfbookmark[0]{目录}{tjtoc}
\noindent{\tjdisplay\fontsize{22}{28}\selectfont 目录}\hfill{\tjmeta <% if grouped %>按关键词 · <% endif %>共 << count >> 题}%
\vspace{6pt}\tjdoublerule\vspace{4pt}
<% for g in groups %>
<% if g.name %>
\tjtockw{<< g.name >>}{<< g.count >>}{<< g.anchor >>}
<% endif %>
<% for e in g.toc %>
% @@note << e.seq >>
\tjtocitem{<< e.no >>}{<< e.source >>}{<< e.excerpt >>}{<< e.label >>}
% @@end
<% endfor %>
<% endfor %>
% @@end
\clearpage

% ============================================================== 正文
<% for g in groups %>
<% if layout != 'compact' and not loop.first %>
\clearpage
<% endif %>
<% if g.name %>
\tjchapter{<< '%02d' % g.index >>}{<< g.name >>}{<< g.count >>}{<< g.anchor >>}
<% else %>
\markboth{全部题目}{全部题目}\pdfbookmark[0]{全部题目}{tjall}\phantomsection\label{<< g.anchor >>}
<% endif %>
<% for e in g['items'] %>
<% if layout == 'single' and not loop.first %>
\clearpage
<% endif %>
% @@note << e.seq >>
\tjhead{<< e.no >>}{<< e.source >>}{<< e.date >>}{No.<< e.no >><% if e.source %> << e.source >><% endif %>}
\phantomsection\label{<< e.label >>}%
<% if e.keywords %>
\tjkeywords{<% for k in e.keywords %><% if g.name and k == g.name %>\tjchipon{<< k >>}<% else %>\tjchip{<< k >>}<% endif %>\hspace{3pt}<% endfor %>}
<% endif %>
\begin{tjproblem}
<% if e.problem %>
<< e.problem >>
<% if e.image %>
\tcblower
\tjimage{<< e.image >>}{<< e.image_width >>}
<% endif %>
<% elif e.image %>
\tjimage{<< e.image >>}{<< e.image_width >>}
<% else %>
{\color{tjink3}（还没有填写题目）}
<% endif %>
\end{tjproblem}
<% for i in e.ideas %>
<% if i.ai %>
\begin{tjaiidea}
<< i.text >>
\end{tjaiidea}
<% else %>
\begin{tjidea}{<< i.label >>}{tjaccent}
<< i.text >>
\end{tjidea}
<% endif %>
<% endfor %>
% @@end
<% endfor %>
<% if g.see %>
\begin{tjsee}
<% for e in g.see %>
% @@note << e.seq >>
\tjseeitem{<< e.no >>}{<< e.source >>}{<< e.excerpt >>}{<< e.label >>}
% @@end
<% endfor %>
\end{tjsee}
<% endif %>
<% endfor %>

\par\vspace{28pt}
\begin{center}\tikz\fill[tjlvl] (0,0) -- (3pt,3.6pt) -- (6pt,0) -- (3pt,-3.6pt) -- cycle;\end{center}
\end{document}
